% !PRINTED 38
\weberpage{38}
\begin{center}
\S\,11. \textbf{\gesperrt{The three kinds of elliptic integrals.}}
\end{center}
Let now
\weq{1}{f(x) = a_0x^4 + a_1x^3 + a_2x^2 + a_3x + a_4}
be an integral function of the third or fourth degree, and let
\weq{2}{\Omega(x) = \int \frac{\Phi(x)\,dx}{\sqrt{f(x)}} = \int R(x)\,dx,}
where for brevity
\weq{3}{R(x) = \frac{\Phi(x)}{\sqrt{f(x)}}}
has been put, be an elliptic integral. In it $x$ denotes a quantity varying without restriction (and complex as well), and $\sqrt{f(x)}$ has, for every value of $x$ for which $f(x)$ does not vanish, two opposite values. By a \gesperrt{point} we understand not a value of $x$ alone, but a pair of values of $x$ and $\sqrt{f(x)}$ belonging together, that is, a value of $x$ with a definite sign of $\sqrt{f(x)}$. If therefore $x_0$ is any value of $x$, then there are two points $x_0$ when $f(x_0)$ is different from zero, but only one when $f(x_0) = 0$. For $x = \infty$ we determine the points by the sign of $\sqrt{f(x)}:x^2$, and so we have two points $x = \infty$ when $f(x)$ is of the fourth degree, and only one when $f(x)$ is of the third.

By known theorems of the theory of functions there exists, if $f(x_0)$ is different from zero, a development by \gesperrt{Taylor's} theorem, convergent in a certain domain:
\weq{4}{R(x) = A_m(x-x_0)^m + A_{m+1}(x-x_0)^{m+1} + A_{m+2}(x-x_0)^{m+2} + \cdots.}
in which $m$ is a positive or negative integer or else zero, while $A_m, A_{m+1}, A_{m+2}, \dots$ denote constants, of which $A_m$ does not vanish. This is a development in ascending powers. In addition, if $f(x)$ is of the fourth (not of the third) degree, there comes the development in \gesperrt{descending} powers
\weq{5}{R(x) = C_m x^{-m-2} + C_{m+1}x^{-m-3} + C_{m+2}x^{-m-4} + \cdots,}
of which we say that it holds in the neighbourhood of an infinitely distant point.
